\section{A fixed-potential Riesz--Talbot counterexample to the buffered
$N=1$ entropy route}
\label{app:riesz-talbot-fixed-counterexample}

The preceding converter restrictions leave one genuinely nonlinear escape:
concentrate the selected angular state on a set of exponentially small area,
apply a translated reflector there, and use reciprocity to refocus the
transmitted state.  This section proves that the escape is real.  It gives a
single bounded zonal scalar potential in $X_3$ for which the selected
$N=1$ entropy diverges along every cofinal positively smoothed, radially
truncated, double-buffered diagonal.  The result disproves the proposed
fixed-potential selected-channel theorem.  It does not determine the spectral
multiplicity of the limiting operator and therefore does not disprove Simon's
conjecture.

\subsection{Translated barriers and the exact Hardy isometry}

\begin{lemma}[Translated-barrier Hardy isometry]
\label{lem:translated-barrier-hardy}
Let $\mathcal H$ be finite dimensional.  Suppose that $A,B\in
\mathcal B(\mathcal H)$ satisfy
\begin{equation}
 A^*A-B^*B=\Id,\qquad A\text{ is invertible},\qquad
 \|BA^{-1}\|<1,
 \label{eq:barrier-flux-data}
\end{equation}
and let $S$ be unitary.  For $|z|<1$ put
\begin{equation}
 F(z)=(A+zSB)^{-1}.
 \label{eq:barrier-hardy-function}
\end{equation}
Then
\begin{equation}
 \frac1{2\pi}\int_0^{2\pi}
 F(e^{i\phi})^*F(e^{i\phi})\dd\phi=\Id,
 \label{eq:barrier-hardy-isometry}
\end{equation}
and consequently, for every $y\ne0$,
\begin{equation}
 \frac1{2\pi}\int_0^{2\pi}
 \log\frac{\|y\|}{\|F(e^{i\phi})y\|}\dd\phi\geq0.
 \label{eq:barrier-hardy-nonnegative-loss}
\end{equation}

Let $P$ be an orthogonal projection fixed by conjugation and let
\begin{equation}
 A=\Id+(\cosh\rho-1)P,\qquad B=(\sinh\rho)P,\qquad \rho>0.
 \label{eq:localized-hyperbolic-barrier}
\end{equation}
If $S\bar y=y$ and, for some $|\lambda|=1$,
\begin{equation}
 \|(\Id-P)y\|+\|\bar y-\lambda y\|
 \leq\epsilon\|y\|,
 \label{eq:localized-phase-real-state}
\end{equation}
then
\begin{equation}
 \frac1{2\pi}\int_0^{2\pi}
 \log\frac{\|y\|}{\|F(e^{i\phi})y\|}\dd\phi
 \geq\log\cosh\rho-C_\rho\epsilon.
 \label{eq:localized-barrier-positive-loss}
\end{equation}
\end{lemma}

\begin{proof}
Set $C=BA^{-1}$ and $T=SC$.  Factoring on the right gives
\begin{equation}
 F(z)=A^{-1}(\Id+zT)^{-1}
 =A^{-1}\sum_{n\geq0}(-zT)^n.
\end{equation}
The flux identity gives
\begin{equation}
 A^{-*}A^{-1}=\Id-C^*C=\Id-T^*T.
\end{equation}
Parseval and telescoping therefore yield
\begin{equation}
 \frac1{2\pi}\int F^*F
 =\sum_{n\geq0}T^{*n}(\Id-T^*T)T^n=\Id.
\end{equation}
Jensen applied to $\log\|F(e^{i\phi})y\|^2$ proves
\eqref{eq:barrier-hardy-nonnegative-loss}.

If $Py=y$ and $\bar y=\lambda y$, then $Sy=\bar\lambda y$ and
\begin{equation}
 F(e^{i\phi})y
 =\frac{y}{\cosh\rho+
   e^{i\phi}\bar\lambda\sinh\rho}.
\end{equation}
The phase mean of the logarithm of the denominator is
$\log\cosh\rho$.  In the approximate case the inverse in
\eqref{eq:barrier-hardy-function} has norm at most $e^\rho$.
The resolvent identity and~\eqref{eq:localized-phase-real-state} compare
$F(e^{i\phi})y$ uniformly with the preceding scalar expression with error
$C_\rho\epsilon\|y\|$.  Both norms lie between fixed multiples of $\|y\|$,
so the logarithm is uniformly Lipschitz and
\eqref{eq:localized-barrier-positive-loss} follows.
\end{proof}

\subsection{Radial realization at summable charged cost}

The angular operators constructed later in this section are realized by one
scalar radial cell as follows.
The order in which the scales are chosen is important.  First fix the finite
incoming angular space, the number of Riesz factors, the desired error, and
the finite integer ratios between their frequencies.  Next choose the weak
slab lengths.  Only then choose a common odd carrier so large that all of the
frequencies are simultaneously far above the inverse radial scale while the
entire cell remains far below the square of the smallest frequency.  This
leaves room for both the forward and reverse Talbot flights.

\begin{proposition}[A scalar reciprocal Riesz cell]
\label{prop:scalar-reciprocal-riesz-cell}
Let $K\Subset(0,\infty)$, let $\mathcal K$ be a finite-dimensional,
conjugation-invariant zonal space containing $e_0$, and fix $m\geq1$,
$\rho>0$, and $\epsilon>0$.  Outside any prescribed radius $R_-$ there is a
bounded, compactly supported, real zonal potential $v_{m,\epsilon}$ with
vanishing spherical mean at every radius such that the following hold.

\begin{enumerate}
 \item At a chosen tuning point $k_\star\in\operatorname{int}K$, its forward
 Riesz train, localized hyperbolic reflector, and reverse train satisfy the
 hypotheses of Proposition~\ref{prop:robust-reciprocal-cascade} on
 $\mathcal K$, with cell error at most $\epsilon$.
 \item If $s=k_\star/k$, the error is at most
 \begin{equation}
   C_{K,\mathcal K}m|s-1|+\epsilon.
   \label{eq:physical-riesz-cell-chromatic-error}
 \end{equation}
 \item Inserting a sufficiently distant free collar before the cell changes
 its transfer, up to error $\epsilon$, only through the scalar circle phase
 $z=e^{i\vartheta(k)}$; the derivative of $\vartheta$ may be made arbitrarily
 large, uniformly on $K$.
 \item Its charged cost obeys
 \begin{equation}
  \int\!\int_{\mathbb S^2}|v_{m,\epsilon}(r,\omega)|^2
       \dd\omega\dd r
  \leq C_\rho e^{-c m}+\epsilon.
  \label{eq:physical-riesz-cell-cost}
 \end{equation}
\end{enumerate}
The supports of any prescribed countable collection of such cells can be
chosen disjoint.
\end{proposition}

\begin{proof}
We give the scale ledger, including the two approximations used repeatedly
below.  On a slab $[R,R+T]$ put $T^{-1}w(\omega)$, where $w$ is a fixed
finite-band real multiplier of degree at most $D$.  Variation of constants,
uniformly for $k\in K$, gives
\begin{align}
 \mathsf T_{T,w}(k)
 &=\exp\!\left(-\frac{i}{2k}w\right)
   +O_K\!\left(T^{-1}+\frac{D^2T}{R^2}+\frac{T}{R}\right),
 \label{eq:long-weak-mask-transmission}\\
 \mathsf R_{T,w}(k)
 &=O_K\!\left(T^{-1}+\frac{D^2T}{R^2}+\frac{T}{R}\right),
 \label{eq:long-weak-mask-reflection}\\
 \int_R^{R+T}\|T^{-1}w\|_2^2\dd r
 &=T^{-1}\|w\|_2^2.
 \label{eq:long-weak-mask-cost}
\end{align}
The first error is the one-dimensional long-weak-slab error, the second comes
from the angular generator on the slab, and the third from replacing $r$ by
$R$.  The estimates also hold after one $k$ derivative.  On a free collar
$[r_0,r_1]$, the outgoing evolution on every fixed angular block is
\begin{equation}
 \exp\!\left[-\frac{iA_3}{2k}
       \left(\frac1{r_0}-\frac1{r_1}\right)\right]
 +O_K(r_0^{-1}),
 \label{eq:free-collar-angular-flight}
\end{equation}
with the error uniform after the block has been fixed.  These are the standard
Volterra and outgoing Hankel asymptotics; their uniformity is exactly the
finite-block uniformity imposed in the statement.

Apply Proposition~\ref{prop:talbot-riesz-concentration} and first record its
finite odd frequency ratios $1=q_1<\cdots<q_m=Q$.  Choose the mask-slab
lengths so large that the sum of~\eqref{eq:long-weak-mask-cost}, except for
the reflector, is below $\epsilon/10$.  Now choose a common odd carrier
$L_\ast$ so large that, with $L_t=L_\ast q_t$ and with $T_{\max}$ the largest
slab length, there is a radial interval satisfying
\begin{equation}
 R_-+\operatorname{length}(\text{cell})\ll R,\qquad
 QL_\ast\sqrt{T_{\max}}\ll R\ll L_\ast^2.
 \label{eq:riesz-cell-scale-separation}
\end{equation}
This interval is nonempty once $L_\ast\gg Q\sqrt{T_{\max}}$.  Thus the
centrifugal error in every weak mask slab tends to zero, while
\eqref{eq:free-collar-angular-flight} can realize both
$\pi/(2L_t^2)$ and $3\pi/(2L_t^2)$ in the forward and reverse trains.  A
successive enlargement of $L_\ast$, $R$, and the slab lengths makes the sum
of all transmission, reflection, spherical-Talbot, and outgoing errors below
$\epsilon$.

At the focus insert a fixed smooth radial bump $Q_0$ multiplying
\begin{equation}
 \chi_m-\langle\chi_m\rangle_{\mathbb S^2}.
 \label{eq:mean-zero-localized-reflector}
\end{equation}
Choose the bump strength so that the induced two-port coefficient on the
concentrated state is~\eqref{eq:localized-hyperbolic-barrier}.  Subtracting
the spherical mean changes only the scalar channel; a compensating
reflectionless scalar phase slab absorbs that change.  Equations
\eqref{eq:riesz-concentration}--\eqref{eq:riesz-small-support-cost} and the
stability of finite-dimensional transfer matrices show that the resulting
cell has error $O(\epsilon+e^{-cm})$ and reflector cost $C_\rho e^{-cm}$.
All masks have zero mean because the frequencies are odd, and
\eqref{eq:mean-zero-localized-reflector} has zero mean by construction.

The reverse train gives~\eqref{eq:riesz-refocusing}; differentiating the
finite product in $s=k_\star/k$ gives
\eqref{eq:physical-riesz-cell-chromatic-error}.  A free collar of length
comparable to its starting radius conjugates the reflection coefficient by
$e^{2ikd}$, while its angular action is
$O((1/r_0-1/r_1)D^2)=o(1)$ on the already fixed block.  Rechoose the common
carrier beyond that collar and absorb this error into $\epsilon$.  This proves
the phase assertion.  Finally choose the next cell beyond the end of the
preceding one.
This proves disjointness and~\eqref{eq:physical-riesz-cell-cost}.
\end{proof}

\subsection{Uniform lacunary placement and the fixed potential}

The radial translations must work for a continuum of energies and for all
later positive smoothings and Galerkin cutoffs.  The required uniformity is a
compact-family form of the Riemann--Lebesgue lemma.

\begin{lemma}[Uniform torus placement]
\label{lem:uniform-torus-placement}
Let $I$ be a compact interval and let
$\mathcal F\Subset L^1(I;C(\mathbb T^q))$.  Given $\delta>0$ and $M>0$,
there are real frequencies $M<n_1<\cdots<n_q$ such that, uniformly for
$F\in\mathcal F$,
\begin{equation}
 \left|\int_I F(k,e^{in_1k},\ldots,e^{in_qk})\dd k
 -\int_I\int_{\mathbb T^q}F(k,z)\dd z\dd k\right|<\delta.
 \label{eq:uniform-torus-placement}
\end{equation}
The frequencies may be chosen successively: once $n_1,\ldots,n_{q-1}$ have
been fixed to meet earlier finite lists of inequalities, $n_q$ can be taken
arbitrarily large without disturbing those inequalities.  The same conclusion
holds simultaneously for finitely many compact families.
\end{lemma}

\begin{proof}
Approximate the compact family by finitely many functions which are
trigonometric polynomials in $z$ with $L^1(I)$ coefficients.  Choose $n_1$
so that every nonconstant $z_1$ Fourier mode is small.  Having chosen
$n_1,\ldots,n_{t-1}$, every Fourier mode involving a nonzero power of $z_t$
has frequency $\ell_t n_t+O(1)$ and tends to zero as $n_t\to\infty$ by the
Riemann--Lebesgue lemma.  Its zero power is the $z_t$ average.  Induct on $t$
and absorb the finite-net errors.  Earlier inequalities do not contain the
new frequency and are therefore unchanged.  Taking the maximum of finitely
many thresholds proves the final assertion.
\end{proof}

\begin{theorem}[Failure of fixed-potential buffered selected entropy]
\label{thm:fixed-potential-riesz-counterexample}
There are a bounded real zonal potential
$V\in X_3$, with
\begin{equation}
 \int_{\mathbb S^2}V(r,\omega)\dd\omega=0
 \quad\text{for a.e. }r,
 \label{eq:fixed-counterexample-zero-mean}
\end{equation}
and a compact interval $K\Subset(0,\infty)$ with the following property.
For every choice
\begin{equation}
 B_n,R_n,L_n\longrightarrow\infty,qquad
 \frac{L_n}{B_n+R_n+1}\longrightarrow\infty,
 \label{eq:arbitrary-buffered-diagonal}
\end{equation}
let
\begin{equation}
 V_n=\mathbf1_{[1,R_n]}\mathsf S_{B_n}V,
 \label{eq:fixed-counterexample-approximants}
\end{equation}
and let $J_n(k)$ be the Jost matrix of
$H_{R_n,L_n}(V_n)$.  Then
\begin{equation}
 \int_K\log^-\|J_n(k)^{-1}e_0\|^2\dd k
 \longrightarrow+\infty.
 \label{eq:fixed-potential-selected-entropy-divergence}
\end{equation}
In particular, the $N=1$ hypothesis
\eqref{eq:buffered-diagonal-PSR} is false even when one is free to choose the
cofinal positive smoothing, radial truncation, and double-buffered Galerkin
diagonal.
\end{theorem}

\begin{proof}
Fix $k_\star>0$ and a small compact interval $K$ about it.  Take
$m_j=\lceil C_0\log(j+1)\rceil$, with $C_0$ large enough that
\begin{equation}
 \sum_j e^{-cm_j}<\infty.
 \label{eq:riesz-cell-cost-summable}
\end{equation}
Use Proposition~\ref{prop:scalar-reciprocal-riesz-cell} with errors
$\epsilon_j$ satisfying $\sum_j\epsilon_j^\beta<\infty$, and put the resulting
cells on successive disjoint shells.  Their amplitudes are uniformly bounded;
by~\eqref{eq:physical-riesz-cell-cost} and
\eqref{eq:riesz-cell-cost-summable},
\begin{equation}
 \|V\|_{X_3}^2
 =\sum_j\|v_{m_j,\epsilon_j}\|_{L^2_{r,\omega}}^2<\infty.
 \label{eq:fixed-riesz-potential-cost}
\end{equation}
The mean-zero and reality assertions hold cell by cell.

Choose the free radial translations successively by
Lemma~\ref{lem:uniform-torus-placement}.  At the $j$th choice the earlier
cells, a finite range of smoothing degrees, all radial cuts through the
current cell, and all Galerkin cutoffs which resolve its finite angular band
form finitely many compact transfer-matrix families.  A diagonal choice with
error at most $2^{-j}$ therefore makes the physical energy average agree with
the independent circle-phase average, uniformly for every later member of
those families.  Increasing the finite range at each step covers every
positive smoothing degree, every radial cut, and every finite cutoff.  The
high-channel Schur estimate
\eqref{eq:buffered-edge-self-energy-vanishes} supplies the remaining uniform
passage from a resolved finite active block to the full double-buffered
Galerkin model.

Group the cells into dyadic blocks
\begin{equation}
 \mathcal B_q=\{j:2^{q-1}<j\leq2^q\}.
\end{equation}
For $N=2^q$, Proposition~\ref{prop:robust-reciprocal-cascade} applies on
\begin{equation}
 |s-1|\leq w_N
 =\frac{c}{N\log(N+1)\log\log(N+2)}.
\end{equation}
Its conditional version, obtained by starting the proof at the left endpoint
of $\mathcal B_q$, gives circle-phase mean loss at least $c2^q-C$.  Because
$s=k_\star/k$ is bi-Lipschitz near $k_\star$, the corresponding energy window
has length comparable to $w_{2^q}$.  On the complement, the conditional phase
mean is nonnegative by Lemma~\ref{lem:translated-barrier-hardy}.  Hence every
resolved dyadic block contributes at least
\begin{equation}
 c2^qw_{2^q}-C2^{-q}
 \geq \frac{c'}{q\log(q+2)}-C2^{-q}
 \label{eq:dyadic-block-entropy-gain}
\end{equation}
to the energy-averaged selected loss.  The uniform placement error is included
in the summable final term.

A smoothing or radial cut may leave one cell only partially resolved.  Its
circle-phase mean cannot decrease the accumulated loss, again by
\eqref{eq:barrier-hardy-nonnegative-loss}; uniform placement bounds its actual
negative contribution by the assigned summable error.  Cofinality of
\eqref{eq:arbitrary-buffered-diagonal} implies that every fixed dyadic block
is eventually fully resolved and remains far below the Galerkin edge.
Therefore~\eqref{eq:dyadic-block-entropy-gain} and
\begin{equation}
 \sum_{q\geq2}\frac1{q\log q}=+\infty
\end{equation}
give divergence of the signed selected loss.

Finally, pointwise for $t>0$ one has $-\log t\leq\log^-t$.  Apply this with
$t=\|J_n(k)^{-1}e_0\|^2$.  The divergent lower bound for the signed selected
loss therefore implies~\eqref{eq:fixed-potential-selected-entropy-divergence}
directly.  This proves the theorem.
\end{proof}

\begin{corollary}[What the counterexample decides]
\label{cor:fixed-potential-route-closed}
The universal endpoint estimates $\mathrm{SCI}_1$ and $\mathrm{HFI}_1$ are
false in every form which implies the fixed selected-channel entropy bound
\eqref{eq:buffered-diagonal-PSR}.  The exact signed centrifugal and angular
heat interpolation identities remain valid, but their proposed uniform
endpoint bounds do not.  Theorem
\ref{thm:fixed-potential-riesz-counterexample} does not disprove Simon's
conjecture: it controls one fixed probe only, while absolutely continuous
spectral density may survive in moving angular directions.  Any continuation
of the conjecture must therefore be basis-free (for example, a moving-probe
or full-operator entropy), or use a mechanism other than selected-channel
entropy.
\end{corollary}

\subsection{Reciprocal phase synchronization}

The earlier obstruction to stacking localized barriers was a possible
exponential amplification of a small non-real tail.  Reciprocal barriers do the
opposite in phase average.  Let $a=\cosh\rho$, $b=\sinh\rho$, and let $H$ be
the symmetric unitary reciprocal phase on an active block.  An ideal refocused
barrier acts by
\begin{equation}
 \Phi_z(H)=(a\Id+zbH)^{-1}(aH+\bar z b\Id),\qquad |z|=1.
 \label{eq:reciprocal-mobius-map}
\end{equation}
For $h,\lambda\in\mathbb T$, direct cross multiplication gives
\begin{equation}
 |\Phi_z(h)-\Phi_z(\lambda)|
 =\frac{|h-\lambda|}
 {|a+zbh|\,|a+zb\lambda|}.
 \label{eq:mobius-chord-identity}
\end{equation}
At a synchronized phase the projective derivative is
\begin{equation}
 L_z(\lambda)=|a+zb\lambda|^{-2},\qquad
 \frac1{2\pi}\int_0^{2\pi}L_{e^{i\phi}}(\lambda)\dd\phi=1.
 \label{eq:mobius-critical-derivative}
\end{equation}
For every $0<\beta<1$, strict Jensen gives
\begin{equation}
 q_\beta:=\frac1{2\pi}\int_0^{2\pi}
 L_{e^{i\phi}}(\lambda)^\beta\dd\phi<1.
 \label{eq:mobius-fractional-contraction}
\end{equation}
The value is independent of $\lambda$.

\begin{proposition}[Robust reciprocal-cascade entropy]
\label{prop:robust-reciprocal-cascade}
Consider $N$ successive real reciprocal cells on nested finite active blocks.
Suppose that, after the forward and reverse converter gauges, the $j$th cell
has transfer coefficients
\begin{equation}
 A_j(s)=a\Id+E_{j,+}(s),\qquad
 B_j(s)=b\Id+E_{j,-}(s)
 \label{eq:near-scalar-cell-coefficients}
\end{equation}
on the active block, with the exact flux relations and
\begin{equation}
 \|E_{j,+}(s)\|+\|E_{j,-}(s)\|
 \leq\eta_j(s),\qquad
 \eta_j(s)\leq C m_j|s-1|+\epsilon_j,
 \label{eq:near-scalar-cell-error}
\end{equation}
where $m_j\leq C_0\log(j+1)$ and
$\sum_j\epsilon_j^\beta$ is sufficiently small.  Starting from the free
constant channel, let $J_N(s,z)$ be the resulting Jost block.  There are
$c_0,C_1>0$ such that, for
\begin{equation}
 |s-1|\leq w_N,\qquad
 w_N=\frac{c_0}{N\log(N+1)\log\log(N+2)},
 \label{eq:dyadic-coherence-window}
\end{equation}
one has
\begin{equation}
 \int_{\mathbb T^N}
 -\log\|J_N(s,z)^{-1}e_0\|\dd z
 \geq c_0N-C_1.
 \label{eq:robust-cascade-linear-loss}
\end{equation}
Every circle carries normalized Haar measure.
\end{proposition}

\begin{proof}
With zero error and $s=1$, the refocused active Jost block has the scalar
recurrence
\begin{equation}
 d_j=a d_{j-1}+z_jb\overline{d_{j-1}},\qquad d_0=1.
 \label{eq:scalar-hyperbolic-recurrence}
\end{equation}
Thus
\begin{equation}
 \int_{\mathbb T}\log|a d+zb\bar d|\dd z
 =\log|d|+\log a,
 \label{eq:scalar-barrier-mean-gain}
\end{equation}
and the exact mean loss is $N\log a$.

For the perturbed cascade put $H_j=J_j^{-1}\overline{J_j}$ in the refocused
active frame.  Its ideal update is~\eqref{eq:reciprocal-mobius-map}; the
operator-norm perturbation is $O(\eta_j)$.  If the spectrum of $H_j$ lies in
an arc $I$, then the ideal image lies in $\Phi_z(I)$ and
\begin{equation}
 |\Phi_z(I)|=\int_I|a+zbe^{i\vartheta}|^{-2}\dd\vartheta.
 \label{eq:mobius-arc-length}
\end{equation}
Its phase mean is $|I|$.  Strict concavity and
\eqref{eq:mobius-fractional-contraction} give, uniformly for arcs of length at
most a fixed $\ell_0<2\pi$,
\begin{equation}
 \mathbb E_z|\Phi_z(I)|^\beta\leq q|I|^\beta,\qquad q<1.
 \label{eq:arc-fractional-contraction}
\end{equation}

There is only one way a short arc can make a macroscopic excursion: the random
repelling point of $\Phi_z$ must fall within $O(|I|+\eta_j)$ of it.
Equivalently, apply the maximal inequality to the nonnegative arc-length
martingale generated by~\eqref{eq:mobius-arc-length}.  An error arc injected
with length $C\eta_j$ has probability at most
$C\eta_j/\ell_0$ of ever reaching $\ell_0$.  Hence the excursion probability
before time $N$ is at most
\begin{equation}
 C\sum_{j\leq N}\eta_j(s)
 \leq \frac{C}{\log\log(N+2)}+C\sum_j\epsilon_j,
 \label{eq:mobius-excursion-probability}
\end{equation}
where~\eqref{eq:dyadic-coherence-window} and
$\sum_{j\leq N}m_j=O(N\log N)$ were used.  On the complement,
\eqref{eq:arc-fractional-contraction} and subadditivity of $t^\beta$ keep the
reciprocal phase in an arbitrarily short arc, except for a set of phase
histories of arbitrarily small measure.

If $H_j$ lies in an arc of length $o(1)$ around $\lambda_j$, then
$H_j\bar y_j=y_j$ implies
\begin{equation}
 \|\bar y_j-\bar\lambda_j y_j\|=o(1)\|y_j\|.
\end{equation}
Lemma~\ref{lem:translated-barrier-hardy} and
\eqref{eq:near-scalar-cell-error} give a conditional phase-mean loss
$\log a-o(1)$ at such a stage.  At the exceptional stages the conditional
mean remains nonnegative by~\eqref{eq:barrier-hardy-nonnegative-loss}.
Choose the summable error budget small, discard finitely many initial stages,
and sum the conditional means.  This proves
\eqref{eq:robust-cascade-linear-loss}.
\end{proof}

For a physical inner Jost matrix $J$, a reflectionless real converter with
transmission $U$, and $x=J^{-1}e_0$, the state at the reflector is $y=Ux$ and
the reciprocal phase is
\begin{equation}
 S=UJ^{-1}\overline J\,U^T.
 \label{eq:riesz-reciprocal-phase}
\end{equation}
The half-line scattering identity makes $S$ unitary and gives the exact
outgoing-square relation
\begin{equation}
 S\bar y=y.
 \label{eq:riesz-outgoing-square}
\end{equation}

\subsection{The zonal quarter-Talbot Riesz factor}

Put $D=-\partial_x^2$ on the circle and fix $0<a<\pi/4$.  Define
\begin{equation}
 g_s(x)=e^{-i\pi sD/2}e^{-ias\cos x}.
 \label{eq:flat-talbot-factor}
\end{equation}
Even and odd Fourier modes have multipliers $1$ and $-i$ at $s=1$, so
\begin{align}
 g_1(x)&=\cos(a\cos x)-\sin(a\cos x)>0,
 \label{eq:positive-talbot-factor}\\
 p(x)&=g_1(x)^2=1-\sin(2a\cos x),\qquad
 \frac1{2\pi}\int_0^{2\pi}p(x)\dd x=1.
 \label{eq:talbot-riesz-density}
\end{align}
Strict Jensen gives the positive tilted entropy
\begin{equation}
 D_p:=\frac1{2\pi}\int_0^{2\pi}p\log p\dd x>0.
 \label{eq:talbot-positive-entropy}
\end{equation}

\begin{proposition}[Zonal spherical Talbot factor and reciprocal inverse]
\label{prop:zonal-talbot-factor}
Let $\mathcal K\subset C^\infty(\mathbb S^2)$ be finite dimensional,
zonal, and invariant under conjugation.  Uniformly for $s$ in a fixed compact
neighbourhood of $1$ and odd integers $L\to\infty$,
\begin{equation}
 e^{-i\pi sA_3/(2L^2)}
   \bigl(e^{-ias\cos(L\theta)}f\bigr)
 =g_s(L\theta)f+\mathcal R_{L,s}f,\qquad
 \|\mathcal R_{L,s}\|_{\mathcal K\to L^2}
 \leq C_{\mathcal K,a}\frac{\sqrt{\log L}}L.
 \label{eq:spherical-talbot-asymptotic}
\end{equation}
A sign-reversed mask following a free flight of scaled duration
$3\pi/(2L^2)$ gives an operator $W_{L,s}$ such that, for the forward factor
$U_{L,s}$ above,
\begin{align}
 \|W_{L,1}U_{L,1}-\Id\|_{\mathcal K\to L^2}
 &\leq C_{\mathcal K,a}\frac{\sqrt{\log L}}L,
 \label{eq:talbot-reciprocal-inverse}\\
 \|W_{L,s}U_{L,s}-W_{L,1}U_{L,1}\|_{\mathcal K\to L^2}
 &\leq C_{\mathcal K,a}|s-1|.
 \label{eq:talbot-chromatic-bound}
\end{align}
\end{proposition}

\begin{proof}
The identity $\cos(L\theta)=T_L(\cos\theta)$ makes the mask globally smooth.
On zonal functions
\begin{equation}
 A_3=-\partial_\theta^2-\cot\theta\,\partial_\theta.
\end{equation}
Compare its scaled evolution with that of $-\partial_\theta^2$ by Duhamel.
For every smooth periodic $H$,
\begin{equation}
 \|\cot\theta\,\partial_\theta H(L\theta)\|_{L^2(\mathbb S^2)}
 \leq C_H L\sqrt{\log L}.
\end{equation}
The evolution time is $O(L^{-2})$, giving
\eqref{eq:spherical-talbot-asymptotic}; derivatives falling on the fixed
envelope cost only $O(L^{-1})$.  In the flat fast variable
$e^{-i3\pi D/2}=e^{i\pi D/2}$.  Thus the sign-reversed mask after the
$3\pi/2$ flight is the inverse factor at $s=1$.  The same Duhamel comparison
proves~\eqref{eq:talbot-reciprocal-inverse}; differentiation in $s$ on the
scaled finite fast-frequency block proves~\eqref{eq:talbot-chromatic-bound}.
\end{proof}

\begin{proposition}[Uniform Riesz concentration on a finite incoming block]
\label{prop:talbot-riesz-concentration}
Fix $0<\alpha<D_p$, a finite-dimensional conjugation-invariant zonal space
$\mathcal K$, an integer $m$, and $\epsilon>0$.  There are odd,
superlacunary frequencies $L_1<\cdots<L_m$, a real finite-band function
$0\leq\chi_m\leq1$, and forward and reverse Talbot products $U_{m,s}$ and
$W_{m,s}$ such that, with
\begin{equation}
 G_m(\theta)=\prod_{t=1}^m g_1(L_t\theta),\qquad
 E_m=\left\{\sum_{t=1}^m\log p(L_t\theta)\geq\alpha m\right\},
 \label{eq:riesz-product-typical-set}
\end{equation}
one has
\begin{align}
 \|U_{m,1}f-G_mf\|_2
 &\leq\epsilon\|f\|_2,
 \label{eq:riesz-forward-approximation}\\
 \|(1-\chi_m)U_{m,1}f\|_2
 &\leq(\epsilon+Ce^{-cm})\|f\|_2,
 \label{eq:riesz-concentration}\\
 \|\chi_m\|_2^2&\leq Ce^{-cm},
 \label{eq:riesz-small-support-cost}\\
 \|W_{m,1}U_{m,1}-\Id\|_{\mathcal K\to L^2}
 &\leq\epsilon,
 \label{eq:riesz-refocusing}\\
 \|W_{m,s}U_{m,s}-W_{m,1}U_{m,1}\|
 +\|U_{m,s}-U_{m,1}\|
 &\leq C_{\mathcal K,a}m|s-1|.
 \label{eq:riesz-chromatic-error}
\end{align}
The last two norms are from $\mathcal K$ to $L^2$.
\end{proposition}

\begin{proof}
Choose the frequencies successively.  Riemann--Lebesgue, uniformly on the unit
sphere of $\mathcal K$, makes every exponential moment used below as close as
desired to the product moment of independent copies of $\log p$.  On $E_m$,
$\prod p(L_t\theta)\geq e^{\alpha m}$, hence $|E_m|\leq Ce^{-\alpha m}$.
Under the tilted density $p(x)\dd x/(2\pi)$, the mean of $\log p$ is
$D_p>\alpha$.  Chernoff therefore gives, uniformly for $f\in\mathcal K$,
\begin{equation}
 \int_{E_m^c}|G_m|^2|f|^2\dd\omega
 \leq Ce^{-cm}\|f\|_2^2.
 \label{eq:riesz-tilted-large-deviation}
\end{equation}
Apply the positive finite-band spherical smoother
\eqref{eq:positive-angular-smoother} to a slightly enlarged typical set.
For sufficiently large smoothing degree it gives $0\leq\chi_m\leq1$,
\eqref{eq:riesz-small-support-cost}, and the same weighted-tail estimate.
Proposition~\ref{prop:zonal-talbot-factor}, with still more separated
frequencies, proves~\eqref{eq:riesz-forward-approximation} and
\eqref{eq:riesz-concentration}.  Apply the reciprocal factors in reverse order
to obtain~\eqref{eq:riesz-refocusing}.  Telescoping the two products and using
\eqref{eq:talbot-chromatic-bound} proves~\eqref{eq:riesz-chromatic-error}.
\end{proof}
